\documentclass[11pt,a4paper]{article}

% ===============================================================
% PACKAGES
% ===============================================================

\usepackage[utf8]{inputenc}
\usepackage[T1]{fontenc}

\usepackage{graphicx}
\usepackage{multirow}
\usepackage{amsmath,amssymb,amsfonts}
\usepackage{amsthm}
\usepackage{mathrsfs}
\usepackage{xcolor}
\usepackage{textcomp}
\usepackage{manyfoot}
\usepackage{booktabs}

\usepackage{algorithm}
\usepackage{algorithmicx}
\usepackage{algpseudocode}

\usepackage{listings}
\usepackage{url}
\usepackage{hyperref}

\usepackage[
    left=1.5cm,
    right=1.5cm,
    top=2cm,
    bottom=2cm
]{geometry}

% ===============================================================
% HYPERREF SETUP
% ===============================================================

\hypersetup{
    colorlinks=true,
    linkcolor=blue,
    citecolor=blue,
    urlcolor=blue
}

% ===============================================================
% THEOREM ENVIRONMENTS
% ===============================================================

\theoremstyle{plain}
\newtheorem{theorem}{Theorem}

\newtheorem{proposition}[theorem]{Proposition}

\theoremstyle{definition}
\newtheorem{example}{Example}
\newtheorem{remark}{Remark}
\newtheorem{definition}{Definition}

% ===============================================================
% DOCUMENT SETTINGS
% ===============================================================

\raggedbottom

% ===============================================================
% TITLE
% ===============================================================

\title{
    \textbf{
    FairLight: City-Wide Emergency Vehicle Route Prioritization
    via Envy-Free Hierarchical Multi-Agent Actor--Critic
    }
}

\author{
    First Author$^{1,2}$\\
    \texttt{iauthor@gmail.com}
    \and
    Second Author$^{2,3}$\\
    \texttt{iiauthor@gmail.com}
    \and
    Third Author$^{1,2}$\\
    \texttt{iiiauthor@gmail.com}
}

\date{}

% ===============================================================
% DOCUMENT
% ===============================================================

\begin{document}

\maketitle

% ===============================================================
% ABSTRACT
% ===============================================================

\begin{abstract}

Urban traffic signal control is a complex sequential decision-making
problem in which decisions made at one intersection can influence
traffic conditions at neighboring and downstream intersections. The
problem becomes increasingly challenging as the size of the road
network grows because a controller must simultaneously achieve
network-level coordination, local responsiveness, and computational
scalability. Reinforcement learning has emerged as a promising approach
for adaptive traffic signal control, while multi-agent reinforcement
learning (MARL) allows control to be distributed among individual
intersections. However, fully decentralized approaches may lack
sufficient global context, whereas centralized approaches become
difficult to scale.

This paper proposes \textbf{FairLight}, a \textbf{Hierarchical
Region-Aware Multi-Agent Reinforcement Learning framework for Fair and
Adaptive Traffic Signal Control}. The framework introduces a regional
abstraction between
network-level coordination and intersection-level decision-making.
The traffic network is partitioned using graph community detection.
Regional traffic states are subsequently processed by a high-level
Meta-Policy using spatial and temporal representations, while low-level
intersection agents use local observations, neighboring information,
and regional guidance to select signal actions.

A central feature of the proposed framework is the explicit treatment
of \textbf{fairness as a traffic-control objective}. Conventional
evaluation based only on average delay, average travel time, or
throughput can hide excessive delays experienced by a smaller group of
vehicles. The proposed framework therefore incorporates
vehicle-level waiting-time distribution, maximum waiting time, tail
waiting time, and Jain's fairness index into optimization and
evaluation. A multi-objective reward function jointly considers traffic
efficiency, fairness, and emergency responsiveness.

The framework also treats \textbf{region construction as part of the
research problem rather than merely a preprocessing step}. DBSCAN,
Louvain, Leiden, and hybrid region construction methods are
investigated to determine how the choice of regional abstraction
affects downstream traffic-control performance.

\noindent\textbf{Keywords:}
Traffic Signal Control,
Multi-Agent Reinforcement Learning,
Emergency Vehicle Routing,
Community Detection,
Louvain Algorithm,
Fairness-Aware Optimization,
Hierarchical Reinforcement Learning

\end{abstract}


% ===============================================================
% INTRODUCTION
% ===============================================================

\section{Introduction}\label{sec1}

Urban traffic signal control is a complex sequential decision-making
problem where decisions at one intersection propagate to neighboring
and downstream intersections. Traffic demand varies dynamically with
time, incidents, weather, and travel patterns, making static or
rule-based controllers insufficient for real-world conditions.
Adaptive systems such as SCOOT introduced traffic-responsive
coordination \cite{scoot}, while reinforcement learning approaches
have since enabled data-driven signal control
\cite{intellight}. Multi-agent reinforcement learning (MARL) further
distributes control across intersections, improving scalability,
though fully decentralized approaches risk optimizing locally at the
expense of network-wide coordination \cite{thousandlights,ma2c}.

Beyond aggregate congestion, traffic signal decisions simultaneously
affect travel time, delay, queue formation, throughput, fuel
consumption, and emissions. Critically, evaluating controllers solely
through aggregate metrics such as average waiting time conceals how
service is distributed among individual vehicles. Consider two
controllers with vehicle waiting times
$W_A = [10, 10, 10, 10, 10, 100]$ and
$W_B = [25, 25, 25, 25, 25, 25]$: Controller~A achieves a lower mean
delay yet inflicts severe starvation on one vehicle. This motivates
fairness as a first-class control objective rather than a post-hoc
evaluation criterion.

A further dimension of real-world uncertainty that existing MARL
frameworks largely overlook is the presence of emergency vehicles
(EMVs). Most traffic signal control approaches, including hierarchical
methods such as HiLight~\cite{hilight} and cooperative approaches such
as CoLight~\cite{colight}, assume homogeneous vehicle populations and
optimize purely for aggregate efficiency. In practice, ambulances,
fire engines, and police vehicles impose stochastic, high-priority
demands on the network that cannot be resolved by ordinary signal
timing. Delayed EMV passage has measurable consequences: studies
report approximately 700 annual fatalities attributable to late
ambulance responses, and an estimated annual mean of 4{,}500
ambulance-involved crashes in the United States
\cite{nellore2016}. EMVLight demonstrated that jointly optimizing EMV
routing and adaptive signal control within a multi-agent actor-critic
framework substantially reduces emergency response times
\cite{emvlight}, but does not address hierarchical city-wide
coordination or fairness. These considerations establish EMV
responsiveness as a necessary component of any realistic traffic
management framework, not an optional extension. At the same time,
providing unrestricted priority to an EMV may excessively delay
ordinary traffic, potentially transferring congestion to other parts
of the network rather than resolving it. EMV prioritization should
therefore not be treated as an isolated intersection-level problem: it
requires coordination across the surrounding traffic network while
maintaining acceptable service for non-emergency vehicles.

Hierarchical reinforcement learning addresses the
scalability--coordination tradeoff by separating long-horizon
strategic decisions from short-horizon local control
\cite{feudalhrl,hirol}. HiLight is particularly relevant here,
employing a Meta-Policy with Transformer-based spatial and LSTM-based
temporal modelling to coordinate decentralized Sub-Policy agents
across large networks~\cite{hilight}. However, its region
partitioning is treated as a preprocessing artifact rather than a
control-relevant variable, and it neither incorporates EMV priority
nor vehicle-level fairness.

This paper proposes \textbf{FairLight: City-Wide Emergency Vehicle
Route Prioritization}, an \textbf{Envy-Free Hierarchical Multi-Agent
Actor-Critic} framework that jointly addresses these limitations.
The framework partitions the traffic network using graph community
detection methods (DBSCAN, Louvain, Leiden, and hybrid variants),
treating region construction as a research variable whose effect on
downstream control performance is explicitly investigated. A
hierarchical coordination layer issues regional guidance to local
intersection agents, while a multi-objective reward function
simultaneously optimizes traffic efficiency, vehicle-level fairness,
and EMV responsiveness. To the best of our
knowledge, no existing framework jointly integrates hierarchical MARL,
topology-aware regional partitioning, envy-free fairness optimization,
and dynamic multi-EMV routing within a unified city-wide architecture.


% ===============================================================
% RELATED WORK
% ===============================================================

\section{Related Work}\label{sec2}

Conventional traffic signal control has primarily relied on analytical
optimization and rule-based methods for determining optimal signal
timings. Webster's signal timing formulation introduced one of the
earliest mathematical models for fixed-cycle traffic signal
optimization, while later adaptive traffic control systems such as
SCOOT and SCATS dynamically adjusted signal timings using detector-based
traffic measurements and predefined control strategies
\cite{webster1958,scoot,scats,opac}. Max-pressure control further
improved network throughput by selecting signal phases according to
queue pressure differences between incoming and outgoing links
\cite{maxpressure}. Although these approaches are computationally
efficient and have been successfully deployed in practice, they rely
heavily on predefined traffic models and handcrafted optimization
rules. Consequently, they struggle to adapt to highly dynamic urban
environments involving unpredictable congestion, heterogeneous
intersections, accidents, and emergency vehicle requests.

To improve adaptability, several optimization and machine learning
techniques have been introduced for traffic signal control.
Evolutionary optimization methods such as Genetic Algorithms optimize
signal timing plans through iterative search, while Fuzzy Logic
Controllers employ expert-defined rules to respond to uncertain
traffic conditions \cite{genetic,fuzzy}. Later studies adopted machine
learning techniques including Support Vector Machines, Artificial
Neural Networks, and Long Short-Term Memory (LSTM) networks to predict
traffic flow and congestion using historical observations
\cite{anntraffic,svmtraffic,lstmtraffic}. Although these methods
improve prediction accuracy and optimization quality, they primarily
estimate future traffic states rather than learning sequential traffic
control policies. Furthermore, they require extensive historical
training data and cannot continuously adapt to rapidly changing
traffic conditions.

Recent advances in deep reinforcement learning have significantly
improved adaptive traffic signal control by enabling traffic
controllers to learn directly through interaction with the traffic
environment. IntelliLight demonstrated the effectiveness of Deep
Q-Networks for adaptive traffic signal control at isolated
intersections \cite{intellight}. PressLight integrated max-pressure
traffic theory into deep reinforcement learning to maximize traffic
throughput while reducing handcrafted reward engineering
\cite{presslight}. FRAP modeled traffic control as a phase competition
problem and exploited traffic phase symmetry to improve learning
efficiency \cite{frap}. AttendLight later introduced an
attention-based architecture capable of generalizing across
heterogeneous intersections with varying lane configurations
\cite{attendlight}. MetaLight focused on transferring learned
policies across unseen intersections \cite{metalight}, while
Advanced-XLight demonstrated that effective traffic state
representation can significantly influence performance in
reinforcement-learning-based traffic signal control
\cite{advancedxlight}. Despite these advances, existing deep
reinforcement learning methods primarily optimize local traffic
efficiency and do not jointly consider emergency vehicle
prioritization, fairness, or large-scale hierarchical coordination.

To improve cooperation among multiple intersections, graph neural
networks and multi-agent reinforcement learning have received
increasing attention. CoLight introduced graph attention networks to
enable communication among neighboring intersections for cooperative
traffic signal control \cite{colight}. MPLight further improved
scalability by introducing parameter sharing across reinforcement
learning agents, allowing efficient deployment over city-scale
transportation networks \cite{mplight}. HiLight extended this concept
by proposing a hierarchical reinforcement learning framework
consisting of Meta-Policy and Sub-Policy agents for coordinating
thousands of intersections efficiently \cite{hilight}. CityFlow
provides a large-scale multi-agent reinforcement learning environment
designed specifically for city traffic scenarios and supports
evaluation of scalable traffic-control algorithms \cite{cityflow}.
Although these methods substantially improve scalability and
inter-agent cooperation, their primary objective remains maximizing
traffic efficiency. They neither address emergency vehicle routing
nor incorporate fairness-aware resource allocation across different
regions of a city.

HiLight's hierarchical Meta-Policy/Sub-Policy formulation
\cite{hilight}, like most of the methods discussed above, operates
under a closed-world assumption: it is trained and evaluated on
recurrent congestion patterns and does not model non-recurrent
disruptions such as emergency vehicle preemption or traffic
incidents. This omission is not merely conceptual. Poorly coordinated
EMV preemption has been shown to increase average delay for
cross-traffic by 18--32 seconds per vehicle per preemption event, with
these localized delays cascading into network-wide congestion
\cite{dtpreemption}. Even well-established preemption strategies carry
a measurable efficiency cost: EMVLight's own evaluation shows that
traditional preemption techniques increase average non-EMV travel
time by approximately 10\% relative to unpreempted baselines
\cite{emvlight}. These results indicate that emergency vehicle
disruption is a first-order cost of operating a real traffic network,
not a negligible edge case that a closed-world controller can safely
ignore. Motivated by this, our framework treats EMVs and accident
events as first-class constraints on the control problem rather than
external events layered on top of it, at the cost of an increase in
average travel time and delay relative to a disruption-free
benchmark---a cost we treat as the price of a tighter approximation to
real-world operating conditions.

Emergency vehicle routing and traffic signal pre-emption have also
received considerable research attention. Conventional emergency
routing methods generally employ shortest-path algorithms such as
Dijkstra's algorithm, A*, and various time-dependent shortest path
formulations to determine optimal routes under static or partially
dynamic traffic conditions
\cite{dijkstraemv,astar,tdsp,dynamicrouting,miprouting}. More recent
studies have introduced reinforcement learning-based routing
strategies capable of adapting to changing traffic environments
\cite{rlrouting}. EMVLight jointly optimizes emergency vehicle routing
and adaptive traffic signal control using multi-agent actor--critic
reinforcement learning, significantly reducing emergency response
times while maintaining acceptable traffic efficiency for ordinary
vehicles \cite{emvlight}. Subsequently, hierarchical multi-agent
actor--critic methods have been proposed to further improve emergency
routing through coordinated traffic signal control
\cite{hierarchicalemv}. However, existing emergency traffic management
approaches generally assume decentralized coordination and do not
explicitly consider fairness, topology-aware regional decomposition,
or simultaneous optimization for multiple emergency vehicles
operating within large urban transportation networks. Visual
sensing-based approaches have further demonstrated that incorporating
real-time distance measurement and vehicle counting into traffic
management can substantially improve EMV scheduling and intersection
clearance \cite{nellore2016}.

Fairness has recently emerged as an important objective in intelligent
transportation systems. FairMARL introduced a fairness-aware
multi-agent reinforcement learning framework based on generalized
social welfare functions to reduce disparities among intersections
while maintaining competitive traffic efficiency \cite{fairmarl}.
Although fairness-aware optimization improves equitable resource
allocation, existing approaches primarily focus on balancing
intersection delays and do not integrate emergency vehicle routing,
hierarchical coordination, or topology-aware traffic management into
a unified framework.

Overall, significant progress has been achieved in analytical
optimization, deep reinforcement learning, graph-based coordination,
emergency vehicle routing, and fairness-aware traffic signal control.
However, these objectives have largely been investigated independently.
Classical optimization methods lack adaptability, deep reinforcement
learning approaches prioritize traffic efficiency over fairness,
graph-based coordination methods do not explicitly support emergency
response, emergency vehicle routing methods overlook hierarchical
city-wide coordination, and fairness-aware methods ignore emergency
traffic management.

To the best of our knowledge, no existing framework jointly integrates
\textbf{hierarchical multi-agent reinforcement learning,
topology-aware regional partitioning, envy-free fairness optimization,
and dynamic multi-emergency vehicle routing} within a unified
city-wide traffic management architecture. The proposed framework,
\textbf{FairLight}, addresses these limitations by introducing an
\textbf{Envy-Free Hierarchical Multi-Agent Actor--Critic} architecture
that simultaneously optimizes traffic efficiency, fairness, emergency
response, and scalable coordination across heterogeneous urban
transportation networks.


% ===============================================================
% MULTI-OBJECTIVE REWARD FUNCTION
% ===============================================================

\section{Multi-Objective Reward Function Design}

An effective traffic signal controller must optimize multiple
objectives simultaneously rather than maximizing traffic throughput
alone. Conventional reinforcement learning approaches primarily
optimize queue lengths or average waiting times, often producing
policies that systematically disadvantage certain intersections,
neglect pedestrian requirements, and fail to prioritize emergency
vehicles adequately. To address these limitations, we design a
multi-objective reward function that jointly optimizes traffic
efficiency, fairness, and emergency vehicle priority.

The overall reward at each decision step is defined as

\begin{equation}
R_{\text{total}}
=
\frac{4}{9}R_{\text{eff}}
+
\frac{2}{9}R_{\text{fair}}
+
\frac{3}{9}R_{\text{emergency}},
\end{equation}

where each component corresponds to a distinct traffic management
objective.

\subsubsection*{Traffic Efficiency Reward}

The efficiency component encourages continuous improvement in network
performance by rewarding reductions in congestion rather than
rewarding only low congestion states. It is defined as

\begin{equation}
R_{\text{eff}}
=
-\beta_1 \Delta Q
-\beta_2 \Delta W
+
L_{\text{avg}},
\end{equation}

where $\Delta Q$ denotes the change in total queue length between
successive decision steps, $\Delta W$ represents the change in
cumulative vehicle waiting time, and $L_{\text{avg}}$ is the average
local throughput reward aggregated across all intersections. The
coefficients $\beta_1$ and $\beta_2$ determine the relative importance
assigned to queue reduction and waiting-time reduction.

Unlike conventional reward formulations based solely on absolute queue
lengths, the use of temporal differences provides denser learning
signals by rewarding measurable improvements between consecutive time
steps. This enables faster convergence while encouraging policies that
continuously reduce congestion.

\subsubsection*{Fairness Reward}

Optimizing only traffic efficiency frequently results in certain
approaches receiving excessively long red phases, leading to
persistent vehicle starvation. To ensure equitable service across the
network, a fairness reward is incorporated:

\begin{equation}
\begin{aligned}
R_{\text{fair}}
={}&
-0.3\,\mathrm{Var}(W)
-0.25\max(W-90,0)
-0.15(1-J)\\
&-0.5\,\mathrm{Envy}
-0.5\max(W_{\text{ped}})
-1.2C_{\text{ped}}.
\end{aligned}
\end{equation}

where $\mathrm{Var}(W)$ is the variance of vehicle waiting times,
$J$ denotes Jain's Fairness Index, $\mathrm{Envy}$ represents the
difference between the best-performing neighboring intersection and
the current intersection, $W_{\text{ped}}$ is pedestrian waiting
time, and $C_{\text{ped}}$ measures pedestrian--vehicle conflicts.

Each penalty addresses a different source of inequity. Waiting-time
variance promotes balanced service among approaches, the starvation
penalty prevents excessively long delays, Jain's Fairness Index
encourages equitable signal allocation, the envy term discourages
local policies that sacrifice one intersection for another, while
pedestrian waiting and conflict penalties ensure that pedestrian
mobility and safety are jointly optimized with vehicular traffic.

\subsubsection*{Emergency Vehicle Reward}

Emergency vehicles require substantially higher priority than regular
traffic because delays directly affect emergency response times.
Their reward component is therefore defined as

\begin{equation}
R_{\text{emergency}}
=
-W_{\text{emergency}},
\end{equation}

where $W_{\text{emergency}}$ is the cumulative waiting time
experienced by emergency vehicles. This formulation directly
penalizes any delay encountered by ambulances, fire trucks, or police
vehicles, encouraging the controller to allocate green phases
immediately upon emergency vehicle detection.

\subsubsection*{Reward Weight Selection}

The weighting coefficients $(\frac{4}{9},\,\frac{2}{9},\,\frac{3}{9})$
reflect the relative importance of each optimization objective.
Traffic efficiency receives the highest weight because congestion
reduction remains the primary objective of urban traffic signal
control. Emergency vehicle priority is assigned the second-highest
weight owing to its critical societal importance. Fairness receives
the remaining weight to prevent pathological policies that
consistently disadvantage particular intersections.

Finally, the combined reward is normalized using an exponential
moving average before being supplied to the reinforcement learning
algorithm. Reward normalization reduces reward variance, improves
numerical stability during training, and accelerates policy
convergence.


% ===============================================================
% FUTURE WORKS
% ===============================================================

\section{Future Works}\label{sec6}

\subsection{Enhanced Emergency Vehicle Routing}

The routing mechanism presented in this work improves on EMVLight's
reactive dynamic Dijkstra heuristic by integrating route planning
into the hierarchical Meta-Policy. However, several directions remain
open. At each Meta-Policy decision step, a candidate route set is
generated from current regional pressure readings and the known
destination, and sub-goals are issued to Sub-Policy agents along
candidate corridors to reduce queue lengths before EMV arrival. This
transforms routing from a two-horizon problem (current link and next
link) into a multi-horizon planning problem aligned with the
hierarchical structure. Future work will extend this to settings
where the destination is unknown at dispatch time---for example, in
mass-casualty events where the nearest available hospital is itself
dynamically determined---and to scenarios with more than two
simultaneously active EMVs whose corridors intersect repeatedly. The
Louvain community structure provides a natural abstraction for
corridor planning at the boundary level: because inter-community
transitions represent sparse bottleneck links, the Meta-Policy can
reason about sequences of community transitions rather than
individual links, substantially reducing the planning horizon while
preserving global routing awareness.

\subsection{Envy-Free Traffic Control}

The envy-free regularization introduced in this work operates within
individual Louvain communities during training. Future work will
investigate adaptive scheduling of $\lambda_e$ across training:
starting with a high efficiency weight and gradually increasing the
fairness weight as the policy matures may accelerate convergence
while preserving equity in the final policy. Additionally, the
definition of envy used here is pairwise and local; extending it to
group envy---where a set of intersections collectively prefers
another group's allocation---would capture more subtle structural
inequities in large networks. The interaction between the envy-free
penalty and the EMV pre-emption reward also warrants further study:
during active EMV passage, pre-emption agents and normal agents
necessarily receive asymmetric rewards, and calibrating
$\lambda_e$ dynamically based on EMV proximity could allow the
system to temporarily suspend fairness constraints along the active
corridor while preserving them network-wide.

\subsection{Accident-Based Traffic Control on Real City Maps}

The accident injection model used in this work assigns a fixed
severity parameter $\kappa$ to each incident and samples accident
locations according to a crash probability distribution derived from
historical report data. Future work will model accidents with
evolving severity: an initial partial blockage ($\kappa < 1$) that
transitions to complete blockage ($\kappa = 1$) as secondary
vehicles arrive and then gradually clears as the incident is
resolved. This temporal evolution of supply disruption is a
qualitatively harder problem for the Meta-Policy, requiring it to
track incident lifecycle stages rather than treating each accident
as a static constraint. Evaluation on the Coimbatore city map will be
extended to multi-district scenarios in which accidents in one part
of the city cascade congestion into adjacent zones, stressing the
inter-community coordination capacity of the Meta-Policy. The
combination of accident-induced rerouting, multiple active EMVs, and
envy-free fairness enforcement represents the most realistic and
demanding test scenario for the complete framework, and constitutes
the primary experimental target for the next phase of this research.


% ===============================================================
% REFERENCES
% ===============================================================

\begin{thebibliography}{99}

% ---------------------------------------------------------------
% Conventional Traffic Signal Control
% ---------------------------------------------------------------

\bibitem{webster1958}
F.~V. Webster,
``Traffic Signal Settings,''
Road Research Technical Paper No.~39,
Road Research Laboratory, London, UK, 1958.

\bibitem{scoot}
P.~B. Hunt, D.~I. Robertson, R.~D. Bretherton, and R.~D. Winton,
``SCOOT---A Traffic Responsive Method of Coordinating Signals,''
TRRL Laboratory Report 1014,
Transport and Road Research Laboratory, 1981.

\bibitem{scats}
A.~G. Sims and K.~W. Dobinson,
``The Sydney Coordinated Adaptive Traffic System---
Principles, Methodology, Algorithms,''
in \textit{Proceedings of the International Conference on Road
Traffic Signalling}, pp.~67--70, 1980.

\bibitem{opac}
N.~H.~Gartner,
``OPAC: A Demand-Responsive Strategy for Traffic Signal Control,''
\textit{Transportation Research Record},
vol.~906, pp.~75--81, 1983.

\bibitem{maxpressure}
P.~Varaiya,
``Max Pressure Control of a Network of Signalized Intersections,''
\textit{Transportation Research Part C: Emerging Technologies},
vol.~36, pp.~177--195, 2013.

% ---------------------------------------------------------------
% Optimization and Machine Learning
% ---------------------------------------------------------------

\bibitem{genetic}
D.~E. Goldberg,
\textit{Genetic Algorithms in Search, Optimization, and Machine
Learning},
Addison-Wesley, 1989.

\bibitem{fuzzy}
L.~A. Zadeh,
``Fuzzy Logic,''
\textit{Computer},
vol.~21, no.~4, pp.~83--93, 1988.

\bibitem{anntraffic}
M.~J. Dougherty and R.~S. Cobbett,
``A Neural Network Approach to the Prediction of Traffic Flow,''
\textit{Transportation Research Part C},
vol.~2, no.~1, pp.~1--17, 1994.

\bibitem{svmtraffic}
V.~N. Vapnik,
\textit{The Nature of Statistical Learning Theory},
Springer, 1995.

\bibitem{lstmtraffic}
S.~Hochreiter and J.~Schmidhuber,
``Long Short-Term Memory,''
\textit{Neural Computation},
vol.~9, no.~8, pp.~1735--1780, 1997.

% ---------------------------------------------------------------
% Deep Reinforcement Learning for Traffic Signal Control
% ---------------------------------------------------------------

\bibitem{intellight}
H.~Wei, G.~Zheng, H.~Yang, and Z.~Li,
``IntelliLight: A Reinforcement Learning Approach for Intelligent
Traffic Light Control,''
in \textit{Proceedings of the 24th ACM SIGKDD International
Conference on Knowledge Discovery \& Data Mining},
pp.~2496--2505, 2018.

\bibitem{presslight}
H.~Wei, C.~Chen, G.~Zheng, K.~Wu, V.~Gayah, K.~Xu, and Z.~Li,
``PressLight: Learning Max Pressure Control to Coordinate Traffic
Signals in Arterial Network,''
in \textit{Proceedings of the 25th ACM SIGKDD International
Conference on Knowledge Discovery \& Data Mining},
pp.~1290--1298, 2019.

\bibitem{frap}
C.~Zheng, Y.~Xie, H.~Wei, and Z.~Li,
``Learning Phase Competition for Traffic Signal Control,''
in \textit{Proceedings of the 28th International Joint Conference on
Artificial Intelligence}, 2019.

\bibitem{attendlight}
A.~Oroojlooy, M.~Nazari, D.~Hajinezhad, and J.~Silva,
``AttendLight: Universal Attention-Based Reinforcement Learning
Model for Traffic Signal Control,''
in \textit{Advances in Neural Information Processing Systems},
vol.~33, 2020.

\bibitem{metalight}
H.~Wei, G.~Zheng, G.~Yao, and Z.~Li,
``MetaLight: Value-Based Meta-Reinforcement Learning for Traffic
Signal Control,''
in \textit{Proceedings of the AAAI Conference on Artificial
Intelligence}, 2019.

\bibitem{advancedxlight}
``Advanced-XLight: Advanced Traffic Signal Control with Deep
Reinforcement Learning,''
2020.

% ---------------------------------------------------------------
% Multi-Agent Reinforcement Learning and Graph Methods
% ---------------------------------------------------------------

\bibitem{colight}
H.~Wei, N.~Xu, H.~Zheng, and Z.~Li,
``CoLight: Learning Network-Level Cooperation for Traffic Signal
Control,''
in \textit{Proceedings of the 28th International Joint Conference on
Artificial Intelligence},
pp.~1405--1411, 2019.

\bibitem{mplight}
H.~Wei, G.~Zheng, H.~Yao, and Z.~Li,
``MPLight: Efficient Multi-Agent Reinforcement Learning for
Large-Scale Traffic Signal Control,''
in \textit{Proceedings of the AAAI Conference on Artificial
Intelligence}, 2020.

\bibitem{ma2c}
T.~Chu, J.~Wang, L.~Codec\`a, and Z.~Li,
``Multi-Agent Deep Reinforcement Learning for Large-scale Traffic
Signal Control,''
\textit{IEEE Transactions on Intelligent Transportation Systems},
2019,
doi:10.1109/TITS.2019.2901791.

\url{https://doi.org/10.1109/TITS.2019.2901791}

\bibitem{thousandlights}
C.~Chen, H.~Wei, N.~Xu, G.~Zheng, M.~Yang, Y.~Xiong, K.~Xu,
and Z.~Li,
``Toward A Thousand Lights: Decentralized Deep Reinforcement
Learning for Large-Scale Traffic Signal Control,''
\textit{Proceedings of the AAAI Conference on Artificial
Intelligence},
vol.~34, no.~04, pp.~3414--3421, 2020,
doi:10.1609/aaai.v34i04.5744.

\url{https://doi.org/10.1609/aaai.v34i04.5744}

\bibitem{cityflow}
H.~Zhang, S.~Feng, C.~Liu, Y.~Ding, Y.~Zhu, Z.~Zhou, W.~Zhang,
Y.~Yu, H.~Jin, and Z.~Li,
``CityFlow: A Multi-Agent Reinforcement Learning Environment for
Large Scale City Traffic Scenario,''
in \textit{Proceedings of the World Wide Web Conference},
pp.~3620--3624, 2019,
doi:10.1145/3308558.3314139.

\url{https://doi.org/10.1145/3308558.3314139}

\bibitem{feudalhrl}
F.~Vezhnevets, S.~M.~Osindero, M.~P.~Schaul, R.~H.~Pascanu,
R.~Heess, M.~Jaderberg, M.~W.~Czarnecki, and M.~S.~Rusu,
``FeUdal Networks for Hierarchical Reinforcement Learning,''
in \textit{Proceedings of the 34th International Conference on
Machine Learning},
pp.~3540--3549, 2017.

\bibitem{hirol}
S.~Nachum, S.~Gu, H.~Lee, and S.~Levine,
``Data-Efficient Hierarchical Reinforcement Learning,''
in \textit{Advances in Neural Information Processing Systems},
vol.~31, 2018.

\bibitem{hilight}
Y.~Zhu, H.~Wen, G.~Min, and M.~Luo,
``HiLight: A Hierarchical Reinforcement Learning Framework with
Global Adversarial Guidance for Large-Scale Traffic Signal
Control,''
\textit{arXiv preprint arXiv:2506.14391}, 2025.

\url{https://arxiv.org/abs/2506.14391}

% ---------------------------------------------------------------
% Emergency Vehicle Routing
% ---------------------------------------------------------------

\bibitem{dijkstraemv}
E.~W. Dijkstra,
``A Note on Two Problems in Connexion with Graphs,''
\textit{Numerische Mathematik},
vol.~1, pp.~269--271, 1959.

\bibitem{astar}
P.~E. Hart, N.~J. Nilsson, and B.~Raphael,
``A Formal Basis for the Heuristic Determination of Minimum Cost
Paths,''
\textit{IEEE Transactions on Systems Science and Cybernetics},
vol.~4, no.~2, pp.~100--107, 1968.

\bibitem{tdsp}
``Time-Dependent Shortest Path Methods for Dynamic Traffic
Routing,''
2020.

\bibitem{dynamicrouting}
``Dynamic Routing for Emergency Vehicles in Urban Transportation
Networks,''
2020.

\bibitem{miprouting}
``Mixed Integer Programming Approaches for Emergency Vehicle
Routing,''
2020.

\bibitem{rlrouting}
``Reinforcement Learning-Based Emergency Vehicle Routing,''
2020.

% ---------------------------------------------------------------
% Emergency Vehicle Routing + Traffic Signal Control
% ---------------------------------------------------------------

\bibitem{emvlight}
H.~Su, Y.~D.~Zhong, J.~Y.~J. Chow, B.~Dey, and L.~Jin,
``EMVLight: A Multi-Agent Reinforcement Learning Framework for an
Emergency Vehicle Decentralized Routing and Traffic Signal Control
System,''
\textit{Transportation Research Part C: Emerging Technologies},
vol.~146, p.~103955, 2023.

\url{https://doi.org/10.1016/j.trc.2022.103955}

\bibitem{dtpreemption}
H.~Su, H.~Deng, and Y.~Sun,
``Emergency Preemption Without Online Exploration: A Decision
Transformer Approach,''
\textit{arXiv preprint arXiv:2603.22315}, 2026.

\url{https://arxiv.org/abs/2603.22315}

\bibitem{hierarchicalemv}
``Hierarchical Multi-Agent Actor--Critic Emergency Routing and
Traffic Signal Control,''
2024.

% ---------------------------------------------------------------
% Fairness
% ---------------------------------------------------------------

\bibitem{fairmarl}
``FairMARL: Fairness-Aware Multi-Agent Reinforcement Learning for
Traffic Signal Control,''
2024.

% ---------------------------------------------------------------
% EMV Visual Sensing
% ---------------------------------------------------------------

\bibitem{nellore2016}
K.~Nellore and G.~P.~Hancke,
``Traffic Management for Emergency Vehicle Priority Based on
Visual Sensing,''
\textit{Sensors},
vol.~16, no.~11, p.~1892, 2016.

\url{https://doi.org/10.3390/s16111892}

% ---------------------------------------------------------------
% Community Detection
% ---------------------------------------------------------------

\bibitem{blondel2008}
V.~D.~Blondel, J.-L.~Guillaume, R.~Lambiotte, and R.~Lefebvre,
``Fast Unfolding of Communities in Large Networks,''
\textit{Journal of Statistical Mechanics: Theory and Experiment},
vol.~2008, no.~10, p.~P10008, 2008.

\bibitem{traag2019}
V.~A.~Traag, L.~Waltman, and N.~J.~van Eck,
``From Louvain to Leiden: Guaranteeing Well-Connected
Communities,''
\textit{Scientific Reports},
vol.~9, no.~1, p.~5233, 2019.

\end{thebibliography}

\end{document}